\documentclass[11pt]{amsart}

\usepackage[T1]{fontenc}
\usepackage{lmodern}
\usepackage{amsmath,amssymb,mathtools}
\usepackage{microtype}
\usepackage[colorlinks=true,linkcolor=blue,citecolor=blue,urlcolor=blue]{hyperref}

\newtheorem{theorem}{Theorem}[section]
\newtheorem{lemma}[theorem]{Lemma}
\newtheorem{corollary}[theorem]{Corollary}
\newtheorem{remark}[theorem]{Remark}

\newcommand{\Sph}{\mathbb S^1}
\newcommand{\vol}{\operatorname{vol}}
\newcommand{\interior}{\operatorname{int}}
\newcommand{\Bspace}{\mathfrak B}
\newcommand{\frust}{\iota}
\newcommand{\Hmeasure}{\mathcal H}

\title[BV compactness and magnetic Cheeger constants]
{BV compactness and positivity of magnetic Cheeger constants}
\date{24 August 2026}

\subjclass[2020]{49J45, 53C23, 58J50}
\keywords{magnetic Steklov operator, magnetic Cheeger constant,
frustration constant, BV compactness, gauge transformation}

\begin{document}

\begin{abstract}
\textbf{Relation to the problem list.}
This manuscript addresses two entries: Question~7 (L7), ``Large-domain
magnetic frustration,'' and Question~6 (L6), ``Magnetic Cheeger positivity.''
It gives a full proof of Question~7, and Question~6 follows as a corollary.
Both conclusions use the same compact-manifold, smooth-boundary, and
nontrivial-gauge hypotheses stated in the list.

Let $M$ be a compact Riemannian manifold with nonempty smooth boundary and let
$\eta$ be a smooth real $1$-form which cannot be gauged away.  We prove the
large-domain frustration conjecture posed by Chakradhar, Gittins, Habib and
Peyerimhoff: domains which exhaust $M$ in volume and whose interior perimeter
tends to zero have frustration bounded away from zero.  The proof is a BV
compactness argument.  A nearly minimizing circle-valued gauge is extended by
a constant across the complement of the domain.  Its total magnetic variation
tends to zero, and the limit is a global smooth gauge, a contradiction.  As a
consequence, both magnetic Cheeger constants are strictly positive whenever
$\eta$ is not gauge-equivalent to zero.
\end{abstract}

\maketitle

\section{Definitions and statement}

Let $(M^m,g)$ be a compact Riemannian manifold with nonempty smooth boundary.
All volumes and hypersurface measures below are computed using $g$.  For a
smooth real $1$-form $\eta$, put
\[
  d^\eta u:=du+i\eta u
\]
for complex-valued functions $u$.  The space of gauge forms is
\[
  \Bspace_M
  :=\left\{\frac{d\tau}{i\tau}:\tau\in C^\infty(M,\Sph)\right\}.
\]
If $D\subset M$ is a smooth relative domain, its magnetic frustration is
\[
  \frust^\eta(D)
  :=\inf_{\tau\in C^\infty(D,\Sph)}
       \int_D |d^\eta\tau|\,d\vol.
\]
We write
\[
  \partial_I D:=\partial D\cap\interior(M),
  \qquad
  \partial_E D:=\partial D\cap\partial M,
\]
and use $|\partial_I D|$ and $|\partial_E D|$ for their
$(m-1)$-dimensional measures.  The magnetic Cheeger constants introduced in
\cite{CGHP} are
\[
 h^\eta(M)
 :=\inf_D\frac{\frust^\eta(D)+|\partial_I D|}{|D|},
 \qquad
 (h^\eta)'(M)
 :=\inf_D\frac{\frust^\eta(D)+|\partial_I D|}{|\partial_E D|},
\]
where the infima run over nonempty smooth relative domains.  A quotient with
zero denominator is understood as $+\infty$.

The conjecture in Section~4.1 of \cite{CGHP} is the following.  We prove it
below.

\smallskip
\noindent\emph{Status note.}
The published 2025 version of \cite{CGHP} states both the large-domain claim
and the resulting positivity question as open.  The argument below is
self-contained; no claim about priority beyond the cited record is intended.

\begin{theorem}[Large-domain magnetic frustration]\label{thm:large-domain}
Suppose that $\eta\notin\Bspace_M$.  There are constants
$\varepsilon,\delta>0$ such that every nonempty smooth relative domain
$D\subset M$ satisfying
\[
 |D|\geq(1-\varepsilon)|M|
 \quad\text{and}\quad
 |\partial_I D|\leq\varepsilon
\]
also satisfies
\[
 \frust^\eta(D)\geq\delta.
\]
\end{theorem}

\section{The compactness lemma}

We regard a complex-valued BV function as an $\mathbb R^2$-valued BV
function.  All BV derivatives on $M$ are \emph{relative} derivatives.  Thus a
jump across $\partial M$ creates no variation measure; only interfaces inside
$M$ contribute.

For $u\in BV(M;\mathbb C)$ define its magnetic derivative measure by
\[
  \mathcal D^\eta u:=Du+i\eta u\,d\vol.
\]

\begin{lemma}[Vanishing magnetic variation produces a global gauge]
\label{lem:compactness}
Let $D_j\subset M$ be smooth relative domains such that
\[
 |M\setminus D_j|\longrightarrow0,
 \qquad
 |\partial_I D_j|\longrightarrow0.
\]
Suppose that there are maps $\tau_j\in C^\infty(D_j,\Sph)$ for which
\[
 \int_{D_j}|d^\eta\tau_j|\,d\vol\longrightarrow0.
\]
Then $\eta\in\Bspace_M$.
\end{lemma}

\begin{proof}
Extend each $\tau_j$ by the fixed value $1\in\Sph$ and set
\[
 u_j:=
 \begin{cases}
   \tau_j&\text{on }D_j,\\
   1&\text{on }M\setminus D_j.
 \end{cases}
\]
The finite-energy condition implies $\tau_j\in W^{1,1}(D_j)$, since
\[
 \int_{D_j}|d\tau_j|\,d\vol
 \leq \int_{D_j}|d^\eta\tau_j|\,d\vol
       +\int_{D_j}|\eta|\,d\vol.
\]
Consequently $\tau_j$ has BV traces on the smooth interface.  The standard
trace and gluing formula (see \cite{AFP}), applied in hypersurface charts,
gives a complex $1$-form measure $J_j$ supported on $\partial_I D_j$ such that
\[
 Du_j=\boldsymbol 1_{D_j}d\tau_j\,d\vol+J_j,
 \qquad
 \|J_j\|(M)\leq2|\partial_I D_j|.
\]
This also covers a two-sided slit: the jump between the two circle-valued
traces has norm at most two.  At a one-sided interface it is the jump between
the trace and the extension value $1$, again of norm at most two.
There is no term on $\partial_E D_j$, because the derivative is relative to
$M$.  Hence
\begin{align}
 \mathcal D^\eta u_j
  ={}&\boldsymbol 1_{D_j}(d\tau_j+i\eta\tau_j)\,d\vol
      +\boldsymbol 1_{M\setminus D_j}i\eta\,d\vol \notag\\
    &+J_j.                                                \label{eq:measure}
\end{align}
It follows that
\begin{equation}
 \|\mathcal D^\eta u_j\|(M)
 \leq
 \int_{D_j}|d^\eta\tau_j|\,d\vol
 +\|\eta\|_{L^\infty(M)}|M\setminus D_j|
 +2|\partial_I D_j|.
 \label{eq:key-estimate}
\end{equation}
The right-hand side tends to zero.

Moreover, $|u_j|=1$ almost everywhere, and
\[
 |Du_j|(M)
 \leq \|\mathcal D^\eta u_j\|(M)+\|\eta\|_{L^1(M)}.
\]
Thus $(u_j)$ is bounded in $BV(M;\mathbb C)$.  By BV compactness on a compact
manifold with boundary, after passing to a subsequence there is
$u\in BV(M;\mathbb C)$ such that
\[
 u_j\longrightarrow u\quad\text{in }L^1(M).
\]
After taking a further subsequence the convergence holds almost everywhere,
so $|u|=1$ almost everywhere.

The $L^1$ convergence implies that $Du_j\to Du$ distributionally and that
$i\eta u_j\,d\vol\to i\eta u\,d\vol$ in total variation.  Passing to the
limit in \eqref{eq:key-estimate} therefore yields
\[
 Du+i\eta u\,d\vol=0.
\]
In particular, the distributional derivative of $u$ is the bounded density
$-i\eta u$.  Hence $u\in W^{1,\infty}(M)$ and
\begin{equation}
 du=-i\eta u \qquad\text{a.e. on }M.                    \label{eq:parallel}
\end{equation}
For completeness, smoothness up to $\partial M$ can be seen directly.  Commute
weak mixed derivatives in \eqref{eq:parallel}; the quadratic terms cancel and
give $(d\eta)u=0$.  Since $|u|=1$, it follows that $d\eta=0$.  On every
contractible interior chart, or boundary half-chart, write $\eta=d\phi$ with
$\phi$ smooth up to the boundary.  Equation \eqref{eq:parallel} then gives
$d(e^{i\phi}u)=0$ distributionally.  Thus $e^{i\phi}u$ is almost everywhere
constant on each connected chart, and $u=ce^{-i\phi}$ there.  Consequently
$u$ has a smooth representative on $M$, and that representative takes values
in $\Sph$ everywhere.

Finally let $\sigma:=\overline u$.  Conjugating \eqref{eq:parallel} gives
$d\sigma=i\eta\sigma$, and therefore
\[
 \eta=\frac{d\sigma}{i\sigma}\in\Bspace_M.
\]
\end{proof}

\section{Proof of the conjecture and the Cheeger consequence}

\begin{proof}[Proof of Theorem~\ref{thm:large-domain}]
Assume the conclusion is false.  For every integer $j\geq2$, apply the
negation with $\varepsilon=j^{-1}$ and $\delta=j^{-1}$.  There is then a
smooth relative domain $D_j$ such that
\[
 |D_j|\geq(1-j^{-1})|M|,
 \qquad
 |\partial_I D_j|\leq j^{-1},
 \qquad
 \frust^\eta(D_j)<j^{-1}.
\]
Choose $\tau_j\in C^\infty(D_j,\Sph)$ within $j^{-1}$ of the infimum in the
definition of frustration.  Then
\[
 \int_{D_j}|d^\eta\tau_j|\,d\vol<2j^{-1}.
\]
All hypotheses of Lemma~\ref{lem:compactness} hold, so
$\eta\in\Bspace_M$, contrary to the assumption.  This proves the theorem.
\end{proof}

The preceding theorem also settles the positivity question posed immediately
after Theorem~2.5 of \cite{CGHP}.

\begin{corollary}[Positivity of the magnetic Cheeger constants]
\label{cor:cheeger}
If $\eta\notin\Bspace_M$, then
\[
 h^\eta(M)>0,
 \qquad
 (h^\eta)'(M)>0,
 \qquad\text{and hence}\qquad
 h^\eta(M)(h^\eta)'(M)>0.
\]
Consequently the Cheeger--Jammes estimate of \cite[Theorem~2.5]{CGHP} is a
strictly positive lower bound:
\[
 \sigma_1^\eta(M)
 \geq \frac18 h^\eta(M)(h^\eta)'(M)>0.
\]
\end{corollary}

\begin{proof}
Choose the constants in Theorem~\ref{thm:large-domain}, shrinking
$\varepsilon$ if necessary so that $0<\varepsilon<1/2$.  Let $h(M)>0$ be the
ordinary relative Cheeger constant.  For any admissible $D$, one of the
following four estimates applies:
\begin{align*}
 |D|\leq\tfrac12|M|
 &\quad\Longrightarrow\quad
 \frac{\frust^\eta(D)+|\partial_I D|}{|D|}\geq h(M),\\
 \tfrac12|M|<|D|<(1-\varepsilon)|M|
 &\quad\Longrightarrow\quad
 \frac{\frust^\eta(D)+|\partial_I D|}{|D|}
 \geq \frac{\varepsilon}{1-\varepsilon}h(M),\\
 |\partial_I D|>\varepsilon
 &\quad\Longrightarrow\quad
 \frac{\frust^\eta(D)+|\partial_I D|}{|D|}
 \geq\frac{\varepsilon}{|M|},\\
 |D|\geq(1-\varepsilon)|M|,
 \quad |\partial_I D|\leq\varepsilon
 &\quad\Longrightarrow\quad
 \frac{\frust^\eta(D)+|\partial_I D|}{|D|}
 \geq\frac{\delta}{|M|}.
\end{align*}
In the second line one applies the ordinary Cheeger inequality to
$M\setminus D$, whose interior boundary agrees with that of $D$.  Taking the
minimum of the four positive constants proves $h^\eta(M)>0$.

Let
\[
 h'(M):=
 \inf_{\substack{D\text{ admissible}\\ |D|\leq |M|/2}}
 \frac{|\partial_I D|}{|\partial_E D|}.
\]
The positivity $h'(M)>0$ is the relative isoperimetric fact used in the
Steklov Cheeger inequality; see \cite[Theorem~1]{Jammes}.  If
$|D|\leq|M|/2$, then
\[
 \frac{\frust^\eta(D)+|\partial_I D|}{|\partial_E D|}
 \geq h'(M).
\]
If $|D|>|M|/2$, then $|\partial_E D|\leq|\partial M|$ and
\[
 \frac{\frust^\eta(D)+|\partial_I D|}{|\partial_E D|}
 \geq h^\eta(M)\frac{|D|}{|\partial_E D|}
 \geq h^\eta(M)\frac{|M|}{2|\partial M|}.
\]
This proves $(h^\eta)'(M)>0$.  The final assertion is exactly the estimate in
\cite[Theorem~2.5]{CGHP}.
\end{proof}

\begin{remark}[Finite-perimeter formulation]
The proof of Lemma~\ref{lem:compactness} uses only the BV trace and gluing
formula.  It therefore remains valid for finite-perimeter sets $D_j$, with
$\partial_I D_j$ replaced by the reduced boundary inside $M$.  Smoothness of
the domains is needed only to match the formulation in \cite{CGHP}.
\end{remark}

\begin{remark}[Where topology enters]
No explicit choice of generators of $H_1(M)$ is needed.  The topological
obstruction is retained in the circle-valued BV limit: if the frustration,
the omitted volume, and the interior perimeter all vanished simultaneously,
the limit equation $du+i\eta u=0$ would itself produce the forbidden global
gauge.  The jump term in \eqref{eq:measure} is precisely the cost of cutting
the topology.
\end{remark}

\begin{thebibliography}{9}

\bibitem{AFP}
L. Ambrosio, N. Fusco and D. Pallara,
\emph{Functions of bounded variation and free discontinuity problems},
Oxford Mathematical Monographs, Oxford University Press, 2000.

\bibitem{CGHP}
T. Chakradhar, K. Gittins, G. Habib and N. Peyerimhoff,
\emph{A note on the magnetic Steklov operator on functions},
Mathematika \textbf{71} (2025), no.~4, e70037,
\href{https://doi.org/10.1112/mtk.70037}{doi:10.1112/mtk.70037};
\href{https://arxiv.org/abs/2410.07462}{arXiv:2410.07462}.

\bibitem{Jammes}
P. Jammes,
\emph{Une in\'egalit\'e de Cheeger pour le spectre de Steklov},
Ann. Inst. Fourier (Grenoble) \textbf{65} (2015), no.~3, 1381--1385.

\end{thebibliography}

\end{document}
